\chapter{Introduction}\label{ch:1}

\section{Motivation}\label{sec:1.1}

A Raven's Progressive Matrix shows eight panels arranged in a three-by-three grid, with the ninth left empty. Each panel holds simple figures, and the figures change along the rows according to rules the solver is not told: a shape stays the same, a size grows by one step, a colour is the sum of the two before it. The solver has to find the rules in the two complete rows, apply them to the third, and pick the missing panel from eight candidates. \citet{raven1938progressive} designed the test to measure reasoning that does not depend on language or schooling, and the rules behind its items were later analysed by \citet{carpenter1990what}.

Machine learning took over the format through generated datasets. In PGM (\citealp{barrett2018measuring}) and RAVEN (\citealp{zhang2019raven}) a program produces every puzzle. For RAVEN it picks one of seven layouts, gives each attribute of the objects (type, size, colour, number and position) one of four rules (Constant, Progression, Arithmetic, Distribute Three), fills the rows accordingly and renders the panels. I-RAVEN (\citealp{hu2021stratified}) repaired the way RAVEN draws its wrong answers, and I-RAVEN-X (\citealp{hersche2024towards}) lengthened the rows to ten panels and widened the value ranges up to a thousand. These three are the datasets of this thesis.

The first neural solvers were black boxes: networks that map the sixteen images of a puzzle to the index of the answer (\citealp{barrett2018measuring}; \citealp{zhuo2020solving}). Two things about them motivate this work.

The first is that they find shortcuts. In RAVEN every wrong candidate differs from the answer in one attribute, so the answer is the candidate that most resembles the others, and a model that never looks at the matrix can find it (\citealp{hu2021stratified}). In our own check, a model that sees only the eight candidates scores 90 to 100\% on RAVEN and is at chance on I-RAVEN. A high score on such a benchmark may therefore say little about reasoning, which is the general problem of shortcut learning (\citealp{geirhos2020shortcut}).

The second is that they need a great deal of data. We trained a vision transformer on I-RAVEN to generate the answer from the image of the matrix. With 300 training puzzles per layout it answers 16\% of unseen puzzles, where guessing gives 12.5\%. With 1,000 it answers 19\%, with 3,000 it answers 46\%, and only with 6,000 does it reach 93\%. The information is in the small sets all the same: a model that is given the objects of each panel, the facet model of Chapter~\ref{ch:4}, reaches 90\% from the same 300 puzzles. It still reaches 90\% when an object recognizer, trained on the labelled panels of those 300 puzzles, reads the images for it, although that recognizer learns from far more labels than the transformer, which sees only the answer.

\citet{chollet2019measure} argues that this is the number that matters. Skill at a fixed task can be bought with unlimited training data or with unlimited prior knowledge, so skill alone says little about a system; what is informative is how much experience it needed, given the priors it started with. We take this as a rule for reporting. Every model in this thesis states what is built into it, and its accuracy is reported against the number of training puzzles.

What can be built in is the subject of a long argument. One side holds that general methods and more computation overtake hand-built structure in the end (\citealp{sutton2019bitter}). The other holds that structure which matches the problem, relational structure in particular, is what lets a model generalise from little data (\citealp{battaglia2018relational}). Geometric deep learning gives the second view a precise form: the structure to build in is a symmetry, a set of transformations of the input under which the answer does not change or changes in a known way (\citealp{bronstein2021geometric}), and there are proofs that a model with the right invariance generalises strictly better (\citealp{elesedy2021provably}) and, for kernel methods, needs fewer examples by up to the size of the group (\citealp{bietti2021sample}; \citealp{tahmasebi2023exact}).

In practice, though, a prior is still chosen by intuition and judged afterwards by the accuracy it gives. That leaves two questions unanswered each time. How many training examples did the prior replace? And what would it have cost had it been the wrong one? A prior that fits its benchmark looks the same, in a results table, as a prior that is right in general.

\section{The question and the approach}\label{sec:1.2}

This thesis asks: \textbf{what is a built-in symmetry worth to a relational model, and can that be known before training?}

Raven's matrices suit the question for a reason that is usually held against them. The data is generated, so its symmetries are known exactly. The two complete rows of a puzzle can be exchanged. In a layout with two groups of objects, the groups can be exchanged. Size and colour follow the same rules and can be exchanged. Most rules survive when every value is shifted by the same amount, and one does not. Some of these transformations turn a valid puzzle into another puzzle the generator could have produced, some into a valid puzzle it never produces, and some into a puzzle whose right answer is a different one. A symmetry can thus be the right one, a harmless one or a wrong one (correct, extrinsic and incorrect equivariance in the terms of \citealp{wang2023general}), and which it is can be read off the data.

The models are lifted relational neural networks (\citealp{sourek2021beyond}), written in PyNeuraLogic. Such a model is a template: a set of weighted rules over variables. For each puzzle the template is grounded on the facts that describe that puzzle, and the result is a neural network whose shape follows the puzzle: more objects give more neurons, but the weights belong to the rules and are the same everywhere. Like a graph neural network (\citealp{zhou2020graph}), the template shares its weights over a structure; unlike one, it lets the designer say which structure, rule by rule. Whatever a rule refers to only through a variable, it treats alike. So in a template a symmetry is a decision about what to name.

Our claim is that the worth of such a decision can be computed. A symmetry merges the inputs it maps onto one another into orbits, and a model that has the symmetry built in must treat all inputs of an orbit alike. If the rule to be learned treats them alike too, the model has only as many things to learn as there are orbits, and the number of orbits measures the data it needs. If the rule asks for something different on each input of an orbit, the model can be right on at most one of them, and the number of orbits over the number of inputs is a ceiling on its accuracy that no amount of data lifts. Both numbers come from counting, before any model is trained.

\section{Contributions}\label{sec:1.3}

\textbf{A method for building symmetries into a relational model.} We prove that the output of a template does not change when constants that it does not name are renamed throughout a puzzle's facts (Chapter~\ref{ch:3}). The statement is the familiar invariance of relational models under isomorphism and its proof is short; what we add is its use as one test that every symmetry in the thesis passes or fails, before any training. A transformation of the puzzle that amounts to such a renaming is then a symmetry of the model. Leaving something unnamed builds in every permutation of it, and background facts narrow that group down: facts that order the values, for instance, leave only the shifts. The symmetry is exact and holds for every setting of the weights, which we confirm on random weights to about 10⁻¹⁶. With this method we build the facet model (Chapter~\ref{ch:4}), a graph for each puzzle with one node per attribute of each object. It generates a description of the missing panel and only then picks the closest candidate, so it is never told the rules and never sees the candidates, and the shortcut above is closed by construction. On the I-RAVEN test set it answers 97.2\% of puzzles from ground-truth object labels, and 97.0\% end to end, when a recognizer built on the same transformer as our black-box baseline names the objects of every panel, the candidates' included (Chapter~\ref{ch:5}). On the same test puzzles the black-box transformer reaches 93.1\%, with six times the weights of the whole pipeline. The two are not given the same supervision: the recognizer learns from the labelled objects of every panel, the transformer from the label of the answer only. Given the same labels, a plain transformer over the recognizer's output answers 19\% from 300 puzzles per layout, where the facet model answers 90\%, and needs about ten times the data to catch up (Section~\ref{sec:5.4}).

\textbf{A price for a symmetry.} The orbit count above gives, for a pair of model and rule, the data a symmetry saves and the ceiling a mismatched one sets (Chapter~\ref{ch:3}). It needs no network: a learner that only counts, one answer per orbit, realises it. It describes a learner that shares nothing across orbits, which the value models of Chapter~\ref{ch:7} do by construction. A network reaches the count only if it can also tell the orbits apart, and Chapter~\ref{ch:7} shows a case where it could not until the part of the model that describes a row was changed. The ceiling is the error bound of \citet[Theorem 4.3]{wang2023general} for a symmetry the task does not have, evaluated on a rule's table; the count of the data saved, and a test in which both numbers are predicted, are ours.

\textbf{Measurements that test the price.} The sharpest is a pair of rules built for the purpose (Chapter~\ref{ch:7}). In a row (a, b, c), sum sets c = a + b and midpoint sets c = (a + b)/2. The two have the same form and tables of nearly the same size, and models with no symmetry on the values find them equally hard. They differ in one property: adding the same number to every panel keeps a midpoint row valid and breaks a sum row. A model with that shift built in, a convolution along the axis of values, should then learn midpoint from few puzzles and be unable to learn sum, and the count says how few and how unable. From 600 puzzles the model reaches 91.0\% on midpoint and 6.2\% on sum, against a ceiling of 7.9\% (means of three seeds, here and below). The ceiling follows from the shift together with the way the model reads the puzzle, from the question row alone: a model that measured from one row to another would not have it (Section~\ref{sec:3.6.3}). With separate weights for every value, and so without the shift, the same template reaches 18.0\% and 19.4\%. Trained on 100 puzzles, the model gets none of the 101 midpoint rows right that fall in an orbit absent from its training puzzles, at any seed, and about three quarters of the others. Its weights per offset make that unavoidable; the count says in advance which rows those are. And when addition itself is supplied as background facts, the same model learns to use it and gets sum from 100 puzzles (98.2\%). Chapter~\ref{ch:7} says how much that gives away: a rule's full group of moves is the rule, up to one constant. A symmetry is thus not good or bad in itself. Its worth belongs to the pair of model and rule.

The same law holds on the benchmark (Chapter~\ref{ch:6}). Removing a symmetry the data has costs most when data is scarce: 61 points for the value shift at 100 puzzles, 14 and 13 points for the symmetries between components and between attributes at 100 and 600 (three seeds), and between 0 and 6 points at 6,000. Counted in training puzzles before the runs, the components' symmetry is worth a factor of two, the size of its group, which follows almost from how the samples are formed, and the attributes' between two and three, which does not. The models without them need that much more data. A symmetry of the task that the data lacks, the order of Arithmetic's operands, still adds 7 points at 6,000 puzzles. A symmetry that maps valid puzzles to puzzles with another answer holds the model near the ceiling the count gives, about half of all values, whether it trains on 600 puzzles or on 6,000. And a model that names what it should not, here the position of a panel, learns a shortcut through it: after the two complete rows are exchanged, its accuracy on puzzles with a Distribute rule falls from 98.3\% to 39.5\%, while the symmetric version is unchanged.

\section{What the thesis does not claim}\label{sec:1.4}

It does not present the best solver of Raven's matrices. The facet model is the demonstration that the method yields a working model, not the result.

It does not claim that PyNeuraLogic makes a model more accurate. A graph neural network written in PyTorch with the same symmetries reaches 96.8\% on the I-RAVEN test set, against the facet model's 97.2\%. The gains reported here belong to the symmetries. What the lifted template adds is that a symmetry is declared, not engineered, and that one proof covers every symmetry built this way.

It says nothing about human intelligence. Every score is measured on puzzles from the same generator as the training puzzles, with a fixed list of attributes and with rules that run along rows only, which the tests people take do not guarantee. The rows-only prior has a price: on I-RAVEN puzzles whose matrix is transposed, so that the rules run down the columns, every facet model falls from about 97\% to 45\%. The priors studied here were chosen with the generator in view, and Chapter~\ref{ch:8} separates those a person would bring to any such test from those that only this family of benchmarks rewards.

The recognizer is trained apart from the reasoning model, on the labelled objects of every panel, where the black-box model is given labels for the answer only. Chapter~\ref{ch:5} states what each model is given.

It does not claim that our black-box baseline is the best vision model. It is a plain vision transformer trained from scratch. At the same amount of data, three published architectures for these puzzles, SCL, PredRNet and DRNet, reach 83 to 91\% when retrained under one protocol, and 95 to 98\% as their authors report them (\citealp{malkinski2025airaven}). They pick among the candidates, where ours generates the answer. Section~\ref{sec:5.1} sets the numbers side by side.

The comparisons with little data in Chapters~\ref{ch:6} and \ref{ch:7} have three seeds; most other numbers come from a single training run, and the text says which.

\section{The assignment}\label{sec:1.5}

The assignment asks for symmetry-aware neural architectures that combine object recognition with inductive priors on the logical composition of the objects, built on graph or relational neural networks, and for an improvement on black-box vision models. Chapters~\ref{ch:3} and \ref{ch:4} give the architectures and the priors, and Section~\ref{sec:4.3} the object recognition. Chapter~\ref{ch:5} gives the comparison with a black-box vision model on the same task, where the gain is largest when training puzzles are few. Chapters~\ref{ch:6} and \ref{ch:7} measure what each prior contributes.

\section{Structure}\label{sec:1.6}

Chapter~\ref{ch:2} introduces Raven's matrices and the three datasets, graph neural networks, lifted relational neural networks, and symmetry as inductive bias. Chapter~\ref{ch:3} gives the theory: the renaming theorem, how background facts choose the group, the three kinds of symmetry, and the orbit count. Chapter~\ref{ch:4} describes the task, the data, the recognizer, the facet model and the baselines. Chapter~\ref{ch:5} reports the solver's accuracy, with labels and end to end, against the baselines. Chapter~\ref{ch:6} measures what a matched symmetry buys and what a broken or wrong one costs. Chapter~\ref{ch:7} is the study of sum and midpoint. Chapter~\ref{ch:8} discusses where the priors come from, what is specific to lifted templates, and the limits of the work.
